\section{Methods}
\label{sec:methods}

\subsection{Instrument and firmware baseline}
The instrument under test is HunStat2 (Section~\ref{sec:instrument-design}).
The defects diagnosed below were found and fixed in
\texttt{AD5941\_\allowbreak 25/\allowbreak src/\allowbreak
electrochemical\_\allowbreak methods/\allowbreak c\_ca.cpp},
\texttt{c\_swv.cpp}, and \texttt{c\_dpv.cpp} during this project's firmware
development; no other files were changed in the process. The validation
data referenced from the companion manuscripts in
Section~\ref{sec:validation} was collected on this same firmware, compiled
with the \texttt{rp2040:rp2040} Arduino core via \texttt{arduino-cli} and
flashed to a physical XIAO RP2040 board, which carries both corrections
forward.

\subsection{Diagnostic methodology: comparison against vendor reference code}
\label{sec:diagmethod}
CA and SWV had been reported, in prior internal testing, to produce flat,
noisy dummy-cell output with no discernible decay or peak, while CV and OCP
-- which execute through different, independently-written code -- worked
normally. Because the AD5941 is characterized, commercial silicon, the
working hypothesis was that the defect was specific to how CA/SWV/DPV
configured and read that silicon, not a property of the chip itself; the
most direct way to test that hypothesis was to compare the affected classes
against Analog Devices' own reference firmware for the same measurement
class, register write for register write, rather than re-deriving correct
behavior from the datasheet alone. \texttt{AD5940\_\allowbreak
ChronoAmperometric.c/.h} (639\,+\,111 lines) was read in full and compared
line by line against \texttt{c\_ca.cpp}; \texttt{AD5940\_\allowbreak
SqrWaveVoltammetry.h}'s configuration struct and the relevant portions of
\texttt{AD5940\_\allowbreak SqrWaveVoltammetry.c} were compared against
\texttt{c\_swv.cpp}; \texttt{AD5940\_Ramp.c} was
additionally cross-checked as the reference this project's own
already-working legacy CV engine is derived from, to confirm the LP-loop
pattern recurred across independent vendor examples rather than being an
artifact of one file~\cite{ad5940_examples_github}. No AD5940 reference
example for DPV exists in Analog Devices' published set (Section~\ref{sec:dpv-no-ref}),
so \texttt{c\_dpv.cpp} was instead compared against the now-corrected
\texttt{c\_ca.cpp}/\texttt{c\_swv.cpp} pattern, under the reasoning that DPV
is structurally a staircase-plus-single-sided-pulse variant of the same
underlying LP-loop measurement, not a different signal path.

\subsection{Defect 1: excitation/sensing loop selection}
\label{sec:rootcause}
\texttt{c\_ca.cpp}, \texttt{c\_swv.cpp}, and \texttt{c\_dpv.cpp} configured
\texttt{HSLoopCfg\_Type} -- HSDAC excitation, HSTIA sensing via
\texttt{ADCMUXP\_HSTIA\_P}/\texttt{ADCMUXN\_HSTIA\_N} -- for what is
fundamentally a DC/staircase potential-step measurement. Every Analog
Devices reference example for this class of method instead configures
\texttt{LPLoopCfg\_Type}: the LPDAC's 12-bit \emph{Vbias} and 6-bit
\emph{Vzero} channels set the working-electrode bias point (Vzero, applied
via the LPTIA's virtual-ground input) and the counter-electrode drive
(Vbias, applied through the low-power potentiostat amplifier, LPPA, in
feedback with the reference electrode), while the LPTIA senses current at
\texttt{ADCMUXP\_LPTIA0\_P}/\texttt{ADCMUXN\_LPTIA0\_N}. The applied
potential is set from the LPDAC codes as
\begin{equation}
V_{\mathrm{zero}} = \SI{200}{\milli\volt} + \mathrm{code}_{6}\cdot\mathrm{LSB}_{6},
\qquad \mathrm{LSB}_{6} = 64\cdot\mathrm{LSB}_{12},
\label{eq:vzero}
\end{equation}
\begin{equation}
V_{\mathrm{bias}} - V_{\mathrm{zero}} = \mathrm{code}_{12}\cdot\mathrm{LSB}_{12},
\qquad \mathrm{LSB}_{12} = \frac{\SI{2200}{\milli\volt}}{4095},
\label{eq:vbias}
\end{equation}
matching Analog Devices' own \texttt{AppCHRONOAMPSeqCfgGen()} formula
exactly; the fix reconfigures all three classes' \texttt{ConfigDCMeasurement()}
to the LP loop following Eqs.~\eqref{eq:vzero}--\eqref{eq:vbias}, with the
existing user-facing TIA-selection command (\texttt{r0}--\texttt{r7})
remapped from the eight-value HSTIA table to eight representative LPTIA
$\Rtia$ codes chosen to keep approximately the same resistance range
(Table~\ref{tab:rtia}).

\begin{table}[h]
\centering
\footnotesize
\setlength{\tabcolsep}{4pt}
\caption{TIA-selection command remapping (\texttt{r0}--\texttt{r7}).}
\label{tab:rtia}
\begin{tabular}{@{}p{0.9cm}p{2.3cm}p{2.6cm}p{1.7cm}@{}}
\toprule
Code & Old (HSTIA) & New (LPTIA) & New value \\
\midrule
0 & \SI{200}{\ohm}   & \texttt{LPTIARTIA\_200R} & \SI{110}{\ohm} \\
1 & \SI{1}{\kilo\ohm}  & \texttt{LPTIARTIA\_1K}   & \SI{1}{\kilo\ohm} \\
2 & \SI{5}{\kilo\ohm}  & \texttt{LPTIARTIA\_4K}   & \SI{4}{\kilo\ohm} \\
3 & \SI{10}{\kilo\ohm} & \texttt{LPTIARTIA\_10K}  & \SI{10}{\kilo\ohm} \\
4 & \SI{20}{\kilo\ohm} & \texttt{LPTIARTIA\_20K}  & \SI{20}{\kilo\ohm} \\
5 & \SI{40}{\kilo\ohm} & \texttt{LPTIARTIA\_40K}  & \SI{40}{\kilo\ohm} \\
6 & \SI{80}{\kilo\ohm} & \texttt{LPTIARTIA\_85K}  & \SI{85}{\kilo\ohm} \\
7 & \SI{160}{\kilo\ohm}& \texttt{LPTIARTIA\_160K} & \SI{160}{\kilo\ohm} \\
\bottomrule
\end{tabular}
\end{table}

The one functional difference from the vendor reference: Analog Devices'
\texttt{Vzero} is a user-configurable field (defaulting to \SI{1100}{\milli\volt}
in their example); this project hardcodes $V_{\mathrm{zero}}=\SI{1100}{\milli\volt}$
as a constant and exposes only the equivalent of their \texttt{SensorBias}
as the runtime-settable parameter, driven by the serial \texttt{1<value>}
command. This was a deliberate simplification appropriate to a
host-driven, live-configurable instrument (Section~\ref{sec:deliberate}),
not part of the fix.

\subsection{Defect 2: ADC zero-point convention}
\label{sec:adcfix}
The AD5941's signed ADC codes are offset-binary around midscale: raw code
$0{\times}8000$ ($32768$) represents \SI{0}{\volt} differential, not raw
code $0{\times}0000$. The pre-existing conversion cast the raw code
directly to \texttt{int16\_t}, which places the zero crossing at
$0{\times}0000$ instead of the chip's actual midscale. Because the true
signal sits near the real midscale, this cast turned small, physically
sensible signals into large, discontinuous, non-physical swings whenever
the raw code crossed the $0{\times}8000$ boundary -- a complete,
self-sufficient explanation for the previously-reported ``flat, noisy''
behavior, independent of Defect 1. The fix computes the signed signal as
$(\mathrm{int32\_t})\,\mathrm{rawCode} - 32768$ and recovers current as
\begin{equation}
I = \frac{\mathrm{rawCode} - 32768}{32768}\cdot\frac{V_{\mathrm{ref}}}{\mathrm{PGA}\cdot\Rtia},
\label{eq:current}
\end{equation}
with $V_{\mathrm{ref}}=\SI{1.82}{\volt}$, matching the AD5941's documented
ADC code convention and Analog Devices' own
\texttt{AppCHRONOAMPCalcCurrent()} reference
implementation~\cite{ad5940_examples_github}.

\subsection{Ancillary timing-budget correction}
A third, smaller defect shared the same discovery path: the polling
timeout guarding the SINC2 ``data ready'' flag ($1000\times\SI{10}{\micro\second}=\SI{10}{\milli\second}$)
was shorter than the $\approx\SI{13.3}{\milli\second}$ the configured
decimation chain (SINC2 OSR\,=\,1333, SINC3 OSR\,=\,4) requires to produce
one genuinely fresh sample, computed via the AD5940 driver's own
\texttt{AD5940\_ClksCalculate()}. Conversions were very likely timing out
before observing a fresh result. The timeout budget was raised to
$2500\times\SI{10}{\micro\second}=\SI{25}{\milli\second}$, with comfortable
margin above the calculated minimum, and each measurement now performs a
full ADC power-cycle around the conversion matching Analog Devices' own
sequence, rather than a fixed, empirically-guessed delay.

\subsection{Dummy-cell validation protocol}
\label{sec:protocol}
Both defects were corrected and then checked against the same traceable
dummy-cell protocol used by the two companion
manuscripts~\cite{clement2026linearity,clement2026architecture}: a Randles
network ($R_s=\SI{560}{\ohm}$, $R_{ct}=\SI{10}{\kilo\ohm}$,
$C_{dl}=\SI{33}{\nano\farad}$) connected across the WE/RE/CE terminals in
place of an electrochemical cell, with the instrument's native ASCII
command protocol driven from a Python (\texttt{pyserial}) script rather
than the graphical client. This paper does not repeat that data collection;
Section~\ref{sec:validation} reports the outcome already measured on the
same firmware in those two manuscripts, which is the correct basis for
confirming this paper's fix is retained on the firmware both of them
validate, rather than re-measuring the same instrument a third time.
